\documentclass[11pt,a4paper]{article}
\usepackage{geometry}
\geometry{margin=0.6in}  % 从1in缩小到0.6in
\usepackage{enumitem}
\usepackage{hyperref}
\usepackage{titlesec}
\usepackage{fontawesome}
\usepackage{setspace}

\setstretch{1.0}  % 从1.15减少到1.0
\titleformat{\section}{\normalsize\bfseries}{}{0em}{}[\titlerule]
\titlespacing{\section}{0pt}{6pt}{4pt}  % 减少section spacing

% 减少列表间距
\setlist{nosep, leftmargin=1em}


\begin{document}

%==================
% Header
%==================
\begin{center}
{\LARGE \textbf{Siyuan Meng}}\\[4pt]
Email: siyuanmeng@umass.edu \quad|\quad 
GitHub: https://github.com/siyuanmengmax \\
ResearchGate: https://www.researchgate.net/profile/Max-Siyuan-Meng?ev=hdr_xprf 
\end{center}

%==================
\section*{Education}
%==================
\begin{itemize}[leftmargin=1.2em]
    \item \textbf{University of Massachusetts Amherst}, USA \hfill 2023.09 -- 2026.08 (Expected)  
    Ph.D. in Civil Engineering (Transportation Engineering), Fully funded  
    GPA: 3.98/4.0  

    \item \textbf{Chang'an University}, China \hfill 2020.09 -- 2023.07  
    M.S. in Transportation Engineering  
    GPA: 3.74/4.0

    \item \textbf{Chang'an University}, China \hfill 2016.09 -- 2020.07  
    B.S. in Roadway, Bridge, and River-Crossing Engineering  
    GPA: 3.76/4.0
\end{itemize}

%==================
\section*{Research Interests}
%==================
Transportation systems optimization, pavement asset management, LiDAR-based traffic monitoring and simulation, computer vision, intelligent transportation systems, machine learning applications in civil infrastructure.

%==================
\section*{Research Experience}
%==================
\begin{itemize}[leftmargin=1.2em]
    \item \textbf{Graduate Research Assistant}, UMass Amherst \hfill 2023 -- Present  
    Lead multiple projects on LiDAR-based object detection/tracking, multi-objective optimization for roadway asset management.
    \item \textbf{Graduate Research Assistant}, Chang'an University \hfill 2020 -- 2023  
    Conducted research on pavement performance modeling, multi-year maintenance planning, airport curbside vehicle detection.
    \item \textbf{Teaching Assistant}, UMass Amherst \hfill 2024  
    Course: Programming for Civil Engineers.
\end{itemize}

%==================
\section*{Publications}
%==================
\subsection*{Journal Articles}
\begin{enumerate}[leftmargin=1.2em]
    \item Hu, A., Bai, Q., Chen, L., \textbf{Meng, S.}, Li, Q., \& Xu, Z. (2022). A review on empirical methods of pavement performance modeling. \textit{Construction and Building Materials}, 342, 127968.
    \item \textbf{Meng, S.}, Bai, Q., Chen, L., \& Hu, A. (2023). Multiobjective optimization method for pavement segment grouping in multiyear network-level planning of maintenance and rehabilitation. \textit{Journal of Infrastructure Systems}, 29(1), 04022047.
    \item Cao, R., \textbf{Meng, S.}, Feng, S., Dong, S., \& Bai, Q. (2024). Network-level pavement maintenance and rehabilitation planning considering uncertainties using chance-constrained programming. \textit{International Journal of Pavement Engineering}, 25(1), 2340001.
    \item Bai, Q., Wu, S., Cao, R., \textbf{Meng, S.}, \& Xu, Z. (2022). Review of optimization methods for civil airport passenger boarding processes [in Chinese]. \textit{Journal of Transportation Engineering}, 22(04), 68-88. \href{https://doi.org/10.19818/j.cnki.1671-1637.2022.04.005}{doi}.
    \item Bai, Q., Shao, Y., \textbf{Meng, S.}, Wang, Y., Chen, X., \& Feng, H. (2022). Video-based detection method for illegal pick-up vehicles at airport departure levels [in Chinese]. \textit{Journal of Chang'an University (Natural Science Edition)}, 42(04), 73-86. \href{https://doi.org/10.19721/j.cnki.1671-8879.2022.04.008}{doi}.
\end{enumerate}

\subsection*{Conference Papers}
\begin{enumerate}[leftmargin=1.2em]
    \item \textbf{Meng, S.}, Shao, Y., Cao, R., \& Wu, S. (2022). Multi-objective optimization of network-level pavement maintenance decisions considering user costs [in Chinese]. In \textit{Proc. World Transport Convention (WTC 2022) – Transportation Planning and Interdisciplinary}, pp. 693–701. \href{https://doi.org/10.26914/c.cnkihy.2022.019338}{doi}.
    \item Fang, Z., \textbf{Meng, S.} (2025). Estimating travel mode choices of residents in Stockholm region based on multinomial logit model and mixed logit model. \textit{Proc. SPIE 13575, International Conference on Smart Transportation and City Engineering (STCE 2024)}, 135753V.
    \item Li, C., \textbf{Meng, S.} (2025). Analysis of expressway maintenance decision based on complete information static game theory. \textit{Proc. SPIE 13575, STCE 2024}, 135751T.
\end{enumerate}

%==================
\section*{Patents and Software}
%==================
\begin{itemize}[leftmargin=1.2em]
    \item Wu, S., \textbf{Meng, S.}, Cao, R., \& Bai, Q. (2025). Automatic detection method for illegal vehicles at airport departure levels based on deep learning: 202111642183.3 [Granted Patent, in Chinese].
    \item \textbf{Meng, S.} (2023). Airport curbside illegal vehicle monitoring system [Software Copyright, in Chinese].
\end{itemize}

%==================
\section*{Presentations}
%==================
\begin{itemize}[leftmargin=1.2em]
    \item \textbf{Meng, S.}, Shao, Y., Cao, R., \& Wu, S. (2022). Multi-objective optimization of network-level pavement maintenance decisions considering user costs. Poster presented at WTC 2022, Wuhan, China.
    \item \textbf{Meng, S.} (2025). Multi-objective Optimization Method for Integrated Roadway Asset Maintenance and Rehabilitation Considering Traffic Impacts. Poster presented at TRB 2025 Annual Meeting, Washington D.C., USA.
\end{itemize}

%==================
\section*{Under Review / In Preparation}
%==================
\textit{Under Review:}  
\begin{itemize}[leftmargin=1.2em]
    \item MulDet3D: Multi-Objective Optimization Based Unsupervised Object Detection for Multiple Roadside LiDARs (IEEE T-ITS).
    \item FRGB3D: Fast Reliability-Weighted Gaussian Background Modeling for Roadside LiDAR Traffic Monitoring (Journal of Computing in Civil Engineering).
    \item Multi-objective Optimization Method for Integrated Roadway Asset Management Considering Traffic Dynamics and Environmental Impacts (TRR).
\end{itemize}

\textit{In Preparation:}  
\begin{itemize}[leftmargin=1.2em]
    \item MulTrack3D: Reliability-Enhanced Self-Correcting Multi-Point Tracking for Multiple Roadside LiDARs (Target: IEEE T-ITS).
    \item Quality-Aware Traffic Safety Assessment for Roadside Multi-LiDAR Systems: A Time-Varying Reliability Framework for Infrastructure-Based Surrogate Safety Measures (Target: Accident Analysis \& Prevention).
    \item Personalized VISTA: Vision–Language Safety Assistant
    for Trustworthy Autonomy (Target: NeuroIPS).
    \item Coordinated Multi-modal Emergency Evacuation of Airport Landside Passengers under Unexpected Large Passenger Flow (Target: Journal of Transportation Systems Part A)
    \item Cross-Modal Feature Alignment for LiDAR-NIR Image Fusion in 3D Object Detection (Target: Transportation Research Part C)
\end{itemize}

%==================
\section*{Honors \& Awards}
%==================
\begin{itemize}[leftmargin=1.2em]
    \item Outstanding Graduate Student, Chang’an University (May 2023)
    \item First Prize Scholarship, Chang’an University (Dec 2021)
    \item Second Prize, English Speech Contest (Northwest China) (May 2021)
    \item Third Prize, National English Competition for College Students (May 2021)
    \item Third Prize, Mathematical Modeling Competition, Chang’an University (May 2021)
    \item Excellent Undergraduate Thesis, Chang’an University (Jun 2020)
\end{itemize}

%==================
\section*{Industry Experience}
%==================
\textbf{Shaanxi Hongsidao Transportation Planning and Design Co., Ltd} \hfill Jul 2019 -- Aug 2019 \\
\begin{itemize}[leftmargin=1.2em]
    \item Participated in rural highway design in Ankang, Shaanxi Province, China.
    \item Completed partial design drawings of a highway bridge in Jiangsu Province using AutoCAD.
\end{itemize}

\textbf{Taibai Campus, Chang’an University} \hfill Mar 2019 -- Apr 2019 \\
\begin{itemize}[leftmargin=1.2em]
    \item Completed horizontal, vertical, and cross-section designs for a Class IV highway project.
    \item Conducted GPS-based highway surveying on campus.
\end{itemize}

%==================
\section*{Skills}
%==================
Python, MATLAB, C++, LaTeX, TensorFlow, PyTorch, OpenCV, Open3D, GIS, Carla, SUMO, RoadRunner, ROS2, SLAM, optimization modeling, statistical analysis, machine learning, LiDAR data processing.

\end{document}
